\chapter{A journey into electrodynamics}
\label{cha:appendix_electrodynamics}

In this appendix, we develop the electromagnetic descriptions used in chapter~\ref{cha:fundamentals_b}.
We first introduce the electromagnetic fields in matter and the constitutive relations.
We then examine the role of causality in the optical response and derive the Kramers--Kronig relations.
The working optical relations remain in chapter~\ref{cha:fundamentals_b}.
We use SI units, as in that chapter.

\newpage
\section{Toward electromagnetic response in matter}
\label{sec:toward_electromagnetic_response_in_matter}

% opening paragraph
Chapter~\ref{cha:fundamentals_a} focused on the ground-state electronic problem.
There, the central question was how electrons arrange themselves under a given external potential.
However, optical properties do not follow from the ground-state alone.
When an electromagnetic field is applied,
it drives the electrons away from their static equilibrium and triggers a series of responses.

Therefore, before introducing the dielectric function, we need a language that describes both the applied field and the feedback of matter to that field.
That language is electrodynamics.
The present section does not aim to review a complete lecture of electrodynamics;
instead, we only extract the ingredients required for the discussion of the dielectric function, the Kramers--Kronig relations, and the optical properties in chapter~\ref{cha:fundamentals_b}~\cite{landau2013electrodynamics,griffiths2023introduction}.
We begin with electromagnetic fields in matter, and then gradually move toward a frequency-dependent response function that can be connected to the electronic structure.

We use SI units for the macroscopic electrodynamic equations.
The dielectric function is dimensionless and denotes the relative dielectric response.
We use the angular frequency $\omega$ with the time dependence $\exp(-\mathrm{i}\,\omega t)$,
and the photon energy is $\hbar\omega$.

\subsection{Maxwell equations in matter}

Let us start by the universal structure of the problem.
When an electromagnetic field enters a medium, it drives charge redistribution and induces currents inside the medium, while these induced sources further modify the field itself.
Therefore, the optical response of matter should be described within the framework of electrodynamics~\cite{landau2013electrodynamics,griffiths2023introduction}.

At this stage, we do not explore how a specific medium responds;
instead, we first identify the general field equations that govern electromagnetic dynamics in matter.
They provide the universal backbone of the problem, while the constitutive relations later encode the medium-specific information.

Firstly, we write the macroscopic Maxwell equations for the electric field $\vec{E}(\vec{r},t)$ and magnetic field $\vec{B}(\vec{r},t)$ as:

\begin{equation}
\left\{\begin{aligned}
    \nabla\cdot\vec{E}(\vec{r},t) &= \frac{\rho(\vec{r},t)}{\varepsilon_0} 
    &&\space \text{Gauss's law for electric field} \\
    \nabla\cdot\vec{B}(\vec{r},t) &= 0 
    &&\space \text{Gauss's law for magnetic field} \\
    \nabla\times\vec{E}(\vec{r},t) &= -\frac{\partial}{\partial t}\vec{B}(\vec{r},t) 
    &&\space \text{Faraday's law} \\
    \nabla\times\vec{B}(\vec{r},t) &= \mu_0\vec{J}(\vec{r},t)+\mu_0\varepsilon_0\frac{\partial}{\partial t}\vec{E}(\vec{r},t)
    &&\space \text{Ampère--Maxwell law}
\end{aligned}\right.
\label{maxwell_equations_with_total_sources}
\end{equation}

here, $\rho(\vec{r},t)$ and $\vec{J}(\vec{r},t)$ denote the total charge density and total current density, respectively.

These equations already imply a local conservation law.
More specifically, we take the divergence of the Ampère--Maxwell law,
and use Gauss's law for the electric field to obtain the continuity equation for charge at the macroscopic level~\cite{landau2013electrodynamics}:

\begin{equation}
    \frac{\partial}{\partial t}\rho(\vec{r},t)+\nabla\cdot\vec{J}(\vec{r},t)=0
\label{macroscopic_charge_continuity_equation}
\end{equation}

here, the conservation principle discussed in section~\ref{sec:time_evolution_quantum_mechanics} of appendix~\ref{cha:appendix_quantum} reappears in electrodynamic language.

Secondly, for media, it proves useful to decompose the sources into free and bound contributions:

\begin{equation}
\begin{aligned}
    \rho(\vec{r},t) &= \rho_\mathrm{free}(\vec{r},t)+\rho_\mathrm{bound}(\vec{r},t) \\
    \vec{J}(\vec{r},t) &= \vec{J}_\mathrm{free}(\vec{r},t)+\vec{J}_\mathrm{bound}(\vec{r},t)
\end{aligned}
\label{free_and_bound_source_decomposition}
\end{equation}

here, the free part describes externally supplied or mobile macroscopic sources,
while the bound part describes the response generated internally by the medium.

For discussing the bound electric response, we introduce the polarization field $\vec{P}(\vec{r},t)$ as:

\begin{equation}
    \rho_\mathrm{bound}(\vec{r},t)\coloneqq -\nabla\cdot\vec{P}(\vec{r},t)
\label{bound_charge_from_polarization}
\end{equation}

Accordingly, we define the electric displacement field as:

\begin{equation}
    \vec{D}(\vec{r},t)\coloneqq \varepsilon_0\vec{E}(\vec{r},t)+\vec{P}(\vec{r},t)
\label{electric_displacement_field_definition}
\end{equation}

For the optical response considered here, we neglect magnetic-permeability effects and take the relative magnetic permeability as $\mu_{\mathrm{r}}(\omega)\approx1$.
The electric response is then the central object, and we keep the Maxwell equations in the form:

\begin{equation}
\left\{\begin{aligned}
    \nabla\cdot\vec{D}(\vec{r},t) &= \rho_\mathrm{free}(\vec{r},t) 
    &&\space \text{Gauss's law for electric displacement} \\
    \nabla\cdot\vec{B}(\vec{r},t) &= 0 
    &&\space \text{Gauss's law for magnetic field} \\
    \nabla\times\vec{E}(\vec{r},t) &= -\frac{\partial}{\partial t}\vec{B}(\vec{r},t) 
    &&\space \text{Faraday's law} \\
    \frac{1}{\mu_0}\nabla\times\vec{B}(\vec{r},t) &= \vec{J}_\mathrm{free}(\vec{r},t)+\frac{\partial}{\partial t}\vec{D}(\vec{r},t)
    &&\space \text{Amp\`ere--Maxwell law in matter}
\end{aligned}\right.
\label{maxwell_equations_in_matter}
\end{equation}

Consequently, we can now characterize the structure of the problem.
The Maxwell equations themselves are universal,
while the medium dependence enters through the induced polarization $\vec{P}(\vec{r},t)$, and therefore through the electric displacement field $\vec{D}(\vec{r},t)$.

In many optical problems, one considers the response to an external electromagnetic perturbation in the absence of free charges and free currents inside the medium.
Under this condition, equation~\eqref{maxwell_equations_in_matter} shows that the central issue is no longer the field equations themselves,
but how the medium develops polarization under the applied electric field.
Once the relation between $\vec{P}(\vec{r},t)$ and $\vec{E}(\vec{r},t)$ is specified, the electromagnetic response becomes well defined.
This is precisely the point where macroscopic electrodynamics meets microscopic electronic structure,
and it also provides the starting point for introducing the dielectric function~\cite{landau2013electrodynamics,lima2022usefulness,bernasconi2011first}.

\subsection{Polarization, electric displacement, and constitutive relations}

We have established the universal field equations in the previous subsection.
However, those equations do not explore how matter responds to the applied field.
To close the problem, we still need a relation that connects the induced polarization to the electric field.
This relation is the constitutive relation, and it is here that the medium-specific information enters~\cite{landau2013electrodynamics, lima2022usefulness}.

Here, the central object is the polarization field $\vec{P}(\vec{r},t)$.
It describes how the medium rearranges its internal charges under the applied electric field.
Once $\vec{P}(\vec{r},t)$ is known, equation~\eqref{electric_displacement_field_definition} immediately gives the electric displacement field $\vec{D}(\vec{r},t)$.
Therefore, instead of treating $\vec{D}(\vec{r},t)$ as an independent quantity, we may regard it as a convenient way to encode the electric response of matter.

Firstly, we start from the most general linear-response form.
For any time $t$, the polarization may depend not only on the electric field at the same point,
but also on fields acting earlier and in a surrounding spatial region.
Therefore, before introducing the simplified form used in the present work,
we write the constitutive relation in its general nonlocal form in space and time~\cite{royo2019first}:

\begin{equation}
    P_\alpha(\vec{r},t) =
    \varepsilon_0\sum_\beta\int\dif^3{r'}
    \int_{-\infty}^{t}\dif{t'}
    \,\chi_{\alpha\beta}(\vec{r},t;\vec{r}\,',t')\,E_\beta(\vec{r}\,',t')
\label{general_nonlocal_constitutive_relation}
\end{equation}

here, $\chi_{\alpha\beta}(\vec{r},t;\vec{r}\,',t')$ describes the contribution of the electric field acting at $(\vec{r}\,',t')$ to the polarization that develops later at $(\vec{r},t)$.
In this sense, the medium does not have to respond only at one point and one instant;
instead, it may spread the response over space and retain memory in time.
At the same time, the upper limit $t$ in equation~\eqref{general_nonlocal_constitutive_relation} keeps the causal order explicit: the perturbation acts first, and the response follows.

For the present work, we focus on the linear response in the long-wavelength regime.
Under this approximation, spatial nonlocality drops out, and the constitutive relation retains only the memory of the earlier perturbation.
Therefore, it reduces to the time-domain constitutive relation:

\begin{equation}
    P_\alpha(\vec{r},t)=
    \varepsilon_0\sum_\beta\int_{-\infty}^{t}\dif{t'}\,\chi_{\alpha\beta}(t-t')\,E_\beta(\vec{r},t')
\label{time_domain_constitutive_relation}
\end{equation}

here, $\chi_{\alpha\beta}(t-t')$ denotes the time-dependent susceptibility tensor.
This expression demonstrates that the polarization does not follow the electric field instantaneously;
instead, the medium retains a memory of the earlier perturbation.

Optical experiments usually probe the response of a medium under an oscillating electromagnetic field.
Therefore, once we enter the linear-response regime, the frequency domain provides the natural language for discussing the optical response~\cite{han2023temperature}.
Next, we transform equation~\eqref{time_domain_constitutive_relation} to the frequency domain, and obtain:

\begin{equation}
    P_\alpha(\vec{r},\omega)=\varepsilon_0\sum_\beta\chi_{\alpha\beta}(\omega)\,E_\beta(\vec{r},\omega)
\label{frequency_domain_polarization_susceptibility_relation}
\end{equation}

here, $\chi_{\alpha\beta}(\omega)$ denotes the frequency-dependent electric susceptibility tensor.
This equation already shows that the susceptibility determines the linear electric response of the medium.

Substituting equation~\eqref{frequency_domain_polarization_susceptibility_relation} into equation~\eqref{electric_displacement_field_definition},
we then express the electric displacement field in terms of the electric susceptibility tensor:

\begin{equation}
\begin{aligned}
    D_\alpha(\vec{r},\omega)
    &=\varepsilon_0 E_\alpha(\vec{r},\omega)+P_\alpha(\vec{r},\omega) \\
    &=\varepsilon_0\sum_\beta\left[\delta_{\alpha\beta}+\chi_{\alpha\beta}(\omega)\right]E_\beta(\vec{r},\omega)
\end{aligned}
\label{electric_displacement_from_susceptibility}
\end{equation}

It now shows that the combination $\delta_{\alpha\beta}+\chi_{\alpha\beta}(\omega)$ naturally enters the linear electric response.

This observation directly motivates the definition of the dielectric tensor $\varepsilon_{\alpha\beta}(\omega)$:

\begin{equation}
    \varepsilon_{\alpha\beta}(\omega)
    \coloneqq\delta_{\alpha\beta}+\chi_{\alpha\beta}(\omega)
\label{dielectric_tensor_definition}
\end{equation}

Then, the frequency-domain constitutive relation takes the compact form:

\begin{equation}
    D_\alpha(\vec{r},\omega)
    = \varepsilon_0\sum_\beta\varepsilon_{\alpha\beta}(\omega)E_\beta(\vec{r},\omega)
\label{frequency_domain_constitutive_relation_with_dielectric_tensor}
\end{equation}

The dielectric tensor $\varepsilon_{\alpha\beta}(\omega)$ thus provides a compact description of the linear electric response.
Starting from the polarization field $\vec{P}(\vec{r},t)$,
we introduced the electric susceptibility tensor $\chi_{\alpha\beta}(\omega)$,
and then arrived at the dielectric tensor $\varepsilon_{\alpha\beta}(\omega)$.
Therefore, this tensor emerges naturally from the constitutive relation that closes Maxwell electrodynamics in matter.

Besides its directional dependence, the dielectric response is also dynamical.
In linear response theory, causality constrains this dynamical character,
which means that the polarization at a given time can only depend on the electric field applied at earlier times.
This causal structure becomes important in section~\ref{sec:causality_derivation_basic},
where we derive the Kramers--Kronig relations from the analytic structure of the response.

\subsection{Anisotropy and the dielectric tensor}

The previous subsection presented the dielectric tensor $\varepsilon_{\alpha\beta}(\omega)$ as a compact description of the linear electric response.
We now examine how symmetry constrains its form~\cite{capper2019optical}.
When the medium responds identically along all directions, the dielectric tensor reduces to a scalar form.
If the response changes with direction, the full tensor form must be retained.

Firstly, for an isotropic medium,
the dielectric tensor $\varepsilon_{\alpha\beta}(\omega)$ reduces to a scalar dielectric function $\varepsilon(\omega)$~\cite{born2013principles}:

\begin{equation}
    \varepsilon_{\alpha\beta}(\omega)
    = \varepsilon(\omega)\delta_{\alpha\beta}
\label{isotropic_dielectric_tensor}
\end{equation}

here, the scalar dielectric function $\varepsilon(\omega)$ already determines the full linear electric response,
since the isotropic medium does not distinguish one Cartesian direction from another.
The tensor structure, therefore, no longer carries additional directional information.

Secondly, crystalline media usually distinguish directions.
In that case, anisotropy survives, and we must retain the full tensor form.
When the crystal symmetry allows a diagonal dielectric tensor in the chosen Cartesian axes,
the dielectric tensor $\varepsilon_{\alpha\beta}(\omega)$ takes the principal-axis form~\cite{nye1985physical}:

\begin{equation}
    \left[\varepsilon_{\alpha\beta}(\omega)\right] = \begin{pmatrix}
        \varepsilon_{xx}(\omega) & 0 & 0 \\
        0 & \varepsilon_{yy}(\omega) & 0 \\
        0 & 0 & \varepsilon_{zz}(\omega)
    \end{pmatrix}
\label{principal_axis_dielectric_tensor}
\end{equation}

here, each diagonal component describes the dielectric response along one principal axis,
and anisotropy enters through the difference among these components.

Here, this directional dependence plays a central role~\cite{slavich2024exploring}.
In chapter~\ref{cha:fundamentals_b}, we compare the optical response within the structural plane and perpendicular to that plane.
When the two in-plane directions are symmetry-equivalent,
we introduce the in-plane dielectric function $\varepsilon_{\varparallel}(\omega)$ and the out-of-plane dielectric function $\varepsilon_{\perp}(\omega)$~\cite{nye1985physical,guo2021complete,matthes2016influence}:

\begin{equation}
    \left[\varepsilon_{\alpha\beta}(\omega)\right] = \begin{pmatrix}
        \varepsilon_{\varparallel}(\omega) & 0 & 0 \\
        0 & \varepsilon_{\varparallel}(\omega) & 0 \\
        0 & 0 & \varepsilon_{\perp}(\omega)
    \end{pmatrix}
\label{inplane_outofplane_dielectric_tensor}
\end{equation}

here, the in-plane dielectric function $\varepsilon_{\varparallel}(\omega)$ describes the response to an electric field within the structural plane,
while the out-of-plane dielectric function $\varepsilon_{\perp}(\omega)$ describes the response to an electric field normal to that plane.
This notation directly prepares the comparison between the in-plane and out-of-plane dielectric spectra and optical properties~\cite{guo2021complete,slavich2024exploring}.

This anisotropy also propagates to other optical properties derived from the dielectric tensor~\cite{sujith2023large,slavich2024exploring}.
Once the dielectric response changes with direction, the refractive index, the extinction coefficient, the reflectivity, and the energy-loss spectrum also inherit that directional dependence.
Therefore, the tensor structure does not only refine the notation.
It also determines how we organize and interpret the later optical results.

At this stage, the dielectric tensor $\varepsilon_{\alpha\beta}(\omega)$ has been introduced as the quantity that encodes the directional linear electric response of the medium.
When anisotropy is present, its tensor form makes this directional dependence explicit.

Once the principal directions of the dielectric tensor $\varepsilon_{\alpha\beta}(\omega)$ are chosen as the working axes, the electric field $\vec{E}(\omega)$ can be decomposed as:

\begin{equation}
    \vec{E}(\omega)=\sum_{\lambda}E_{\lambda}(\omega)\,\hat{e}_{\lambda}
\label{electric_field_decomposition_principal_directions}
\end{equation}

where $\hat{e}_{\lambda}$ denotes the unit vector along the principal direction $\lambda$.
In the same basis, the induced polarization $\vec{P}(\omega)$ takes the form:

\begin{equation}
    \vec{P}(\omega)=\varepsilon_0\sum_{\lambda}\chi_{\lambda}(\omega)\,E_{\lambda}(\omega)\,\hat{e}_{\lambda}
\label{polarization_decomposition_principal_directions}
\end{equation}

Accordingly, we denote the dielectric response along each principal direction by $\varepsilon_{\lambda}(\omega)$ and express:

\begin{equation}
    \varepsilon_{\lambda}(\omega)=1+\chi_{\lambda}(\omega)
\label{direction_resolved_dielectric_response}
\end{equation}

This direction-resolved form provides the natural bridge from the tensor description to the dielectric function discussed in section~\ref{sec:from_electronic_response_to_dielectric_function}.

\newpage
\section{Causality and the dielectric response}
\label{sec:causality_derivation_basic}

Section~\ref{sec:causality_and_kramers_kronig_relations} uses causality to connect the real and imaginary parts of the dielectric function.
Here, we establish the analytic structure of the response and derive the full-frequency Kramers--Kronig relations.

\subsection{Causality and analytic structure of response functions}

Firstly, to express causality explicitly, 
we start from the time-domain linear relation between
the polarization field $P_{\alpha}(\vec{r},t)$ and
the electric field $E_{\beta}(\vec{r}\,',t')$~\cite{dressel2002electrodynamics}:

\begin{equation}
    P_{\alpha}(\vec{r},t) = 
    \varepsilon_0\sum_{\beta}\int_{-\infty}^{\infty}\dif{t'}
    \int\dif^3{r'}\,\chi_{\alpha\beta}(\vec{r},\vec{r}\,',t-t')
    E_{\beta}(\vec{r}\,',t')
\label{general_linear_response_polarization}
\end{equation}

here, $\chi_{\alpha\beta}(\vec{r},\vec{r}\,',t-t')$ denotes the time-domain electric susceptibility tensor.
This equation characterizes the polarization response at position $\vec{r}$ and time $t$ to an electric field applied at position $\vec{r}\,'$ and time $t'$.

Therefore, causality requires that the medium responds only to earlier external driving fields~\cite{stefanski2022analytical}:

\begin{equation}
    \chi_{\alpha\beta}(\vec{r},\vec{r}\,',t-t')=0
    \qquad \text{for $t<t'$}
\label{causal_condition_susceptibility_kernel}
\end{equation}

This condition defines a retarded response kernel~\cite{landau2013electrodynamics}.
It excludes any polarization response before the external electric field acts.

In the present work,
we consider the optical long-wavelength limit,
where the spatial nonlocality of the dielectric response becomes negligible~\cite{agranovich2013crystal}.
Afterwards, the linear response then reduces to a local frequency-dependent tensor relation~\cite{franta2023dispersion}:

\begin{equation}
    P_{\alpha}(\omega) =
    \varepsilon_0\sum_{\beta}\chi_{\alpha\beta}(\omega)
    \,E_{\beta}(\omega)
\label{frequency_domain_polarization_response}
\end{equation}

here, $\chi_{\alpha\beta}(\omega)$ denotes the frequency-domain electric susceptibility tensor.
It characterizes the response amplitude for each Cartesian component at angular frequency $\omega$.
Meanwhile, it also describes the phase delay of the response at the same frequency~\cite{koutserimpas2024time}.

We then express the dielectric tensor $\varepsilon_{\alpha\beta}(\omega)$ via the electric susceptibility tensor $\chi_{\alpha\beta}(\omega)$~\cite{franta2020symmetry}:

\begin{equation}
    \varepsilon_{\alpha\beta}(\omega) =
    \delta_{\alpha\beta}
    + \chi_{\alpha\beta}(\omega)
\label{dielectric_tensor_from_susceptibility}
\end{equation}

Subsequently, the causal condition in equation~\eqref{causal_condition_susceptibility_kernel}
gives the frequency-domain susceptibility in the form:

\begin{equation}
\chi_{\alpha\beta}(\omega) =
    \int_{0}^{\infty}
    \chi_{\alpha\beta}(t)
    \,\mathrm{e}^{\,\mathrm{i}\omega t}
    \dif{t}
\label{frequency_domain_susceptibility_from_causality}
\end{equation}

This representation reveals the analytic structure directly.
Since the time integral starts from $t=0$, the response function remains analytic in the upper half of the complex frequency plane~\cite{nussenzveig1972causality}.

For a real time-domain response kernel, the susceptibility $\chi_{\alpha\beta}(\omega)$
and dielectric tensor $\varepsilon_{\alpha\beta}(\omega)$ further satisfy the following condition:

\begin{equation}
\left\{\begin{aligned}
    \chi_{\alpha\beta}(-\omega)
    &= \chi_{\alpha\beta}^{*}(\omega) \\
    \varepsilon_{\alpha\beta}(-\omega)
    &= \varepsilon_{\alpha\beta}^{*}(\omega)
\end{aligned}\right.
\label{reality_condition_response_functions}
\end{equation}

Therefore, in the principal-axis representation,
for each complex dielectric component $\varepsilon_{\lambda\lambda}(\omega)$,
the real part $\varepsilon_{1,\lambda\lambda}(\omega)$ is an even function of frequency,
whereas the imaginary part $\varepsilon_{2,\lambda\lambda}(\omega)$ is an odd function:

\begin{equation}
\left\{\begin{aligned}
    \varepsilon_{1,\lambda\lambda}(-\omega)
    &= \varepsilon_{1,\lambda\lambda}(\omega) \\
    \varepsilon_{2,\lambda\lambda}(-\omega)
    &= -\varepsilon_{2,\lambda\lambda}(\omega)
\end{aligned}\right.
\label{parity_of_dielectric_components}
\end{equation}

At this stage, the previous discussion has established the analytic structure of the dielectric response in frequency space~\cite{nussenzveig1972causality,franta2020symmetry,leon2026momentum}.
We next derive the Kramers--Kronig relations between its real and imaginary parts.

\subsection{Kramers--Kronig relations for the dielectric function}

The previous subsection established the analytic structure of the dielectric response in frequency space.
We now derive the Kramers--Kronig relations for the dielectric function~\cite{toll1956causality}.
Firstly, the causal structure in equation~\eqref{frequency_domain_susceptibility_from_causality} and the reality condition in equation~\eqref{reality_condition_response_functions}
establish a Hilbert-transform relation between the real and imaginary parts of the dielectric response~\cite{stefanski2022analytical,koutserimpas2024time,franta2020symmetry}.
Here, we consider the regular electronic response without a zero-frequency conductivity pole,
with $\varepsilon_{\alpha\beta}(\omega)\to\delta_{\alpha\beta}$ at high frequency.
For a metal, the conductive term must be separated before applying these relations in the form below.
The dielectric function in this derivation then denotes the remaining regular response.
We represent the Kramers--Kronig relations in the frequency-domain form as:

\begin{equation}
\left\{\begin{aligned}
    \varepsilon_{1,\alpha\beta}(\omega)-\delta_{\alpha\beta} &=
    \frac{1}{\mathrm{\pi}}\mathcal{P}
    \int_{-\infty}^{\infty}\dif{\omega'}\,
    \frac{\varepsilon_{2,\alpha\beta}(\omega')}{\omega'-\omega} \\
\varepsilon_{2,\alpha\beta}(\omega) &=
    -\frac{1}{\mathrm{\pi}}\mathcal{P}
    \int_{-\infty}^{\infty}\dif{\omega'}\,
    \frac{\varepsilon_{1,\alpha\beta}(\omega')-\delta_{\alpha\beta}}{\omega'-\omega}
\end{aligned}\right.
\label{kk_general_dielectric_tensor}
\end{equation}

here, $\mathcal{P}$ denotes the Cauchy principal value.
It defines the integral through a symmetric limiting procedure around $\omega'=\omega$~\cite{stefanski2022analytical,koutserimpas2024time}.

This expression identifies the physical roles of the two parts of the dielectric response.
The imaginary part $\varepsilon_{2,\alpha\beta}(\omega)$ characterizes the absorptive channel, 
whereas the real part $\varepsilon_{1,\alpha\beta}(\omega)$ characterizes the dispersive channel.
Consequently, causality unifies these two channels into a single complex response function.

The symmetry relations in equation~\eqref{parity_of_dielectric_components} then reduce the full-frequency integrals to the positive-frequency form in equation~\eqref{kk_positive_frequency_dielectric_component}.
The corresponding static limit and its physical interpretation are discussed in section~\ref{sec:causality_and_kramers_kronig_relations}.
